\documentclass{article}
\usepackage[margin=1in]{geometry}
\usepackage{parskip}
\usepackage{amsmath}
\usepackage{booktabs}
\usepackage{needspace}

\title{ISE 555 HW 5}
\author{Luke Pepin}
\date{October 21 2026}

\begin{document}
\maketitle

Solution sheet: \texttt{HW5.xlsx} (tableaus for Questions 1 and 4). Pivots are boxed, and the row 0 RHS is $-z$.

\section*{Question 1}
The first constraint is $\ge$, so it gets a surplus $e_1$ and an artificial $a_1$. The second is an equality, so it gets an artificial $a_2$.

Minimize: $z = -4x_1 - 2x_2 + Ma_1 + Ma_2$

Subject to:
\begin{itemize}
    \item $3x_1 - 2x_2 - e_1 + a_1 = 4$
    \item $x_1 - 2x_2 + a_2 = 2$
    \item $x_1, x_2, e_1, a_1, a_2 \ge 0$
\end{itemize}

$a_1$ and $a_2$ start in the basis, so row 0 is priced out by subtracting $M$(row 1) and $M$(row 2) from the cost row.

\needspace{9\baselineskip}\textbf{Iteration 0}

\begin{center}
\begin{tabular}{lrrrrrrr}
\toprule
 & $x_1$ & $x_2$ & $e_1$ & $a_1$ & $a_2$ & RHS & ratio \\
\midrule
row 0 & $-4-4M$ & $-2+4M$ & $M$ & $0$ & $0$ & $-6M$ &  \\
row 1 ($a_1$) & $\boxed{3}$ & $-2$ & $-1$ & $1$ & $0$ & $4$ & $4/3$ \\
row 2 ($a_2$) & $1$ & $-2$ & $0$ & $0$ & $1$ & $2$ & $2$ \\
\bottomrule
\end{tabular}
\end{center}

$x_1$ enters, $a_1$ leaves.

\needspace{9\baselineskip}\textbf{Iteration 1}

\begin{center}
\begin{tabular}{lrrrrrrr}
\toprule
 & $x_1$ & $x_2$ & $e_1$ & $a_1$ & $a_2$ & RHS & ratio \\
\midrule
row 0 & $0$ & $-14/3+4M/3$ & $-4/3-M/3$ & $4/3+4M/3$ & $0$ & $16/3-2M/3$ &  \\
row 1 ($x_1$) & $1$ & $-2/3$ & $-1/3$ & $1/3$ & $0$ & $4/3$ & -- \\
row 2 ($a_2$) & $0$ & $-4/3$ & $\boxed{1/3}$ & $-1/3$ & $1$ & $2/3$ & $2$ \\
\bottomrule
\end{tabular}
\end{center}

$e_1$ enters, $a_2$ leaves.

\needspace{9\baselineskip}\textbf{Iteration 2}

\begin{center}
\begin{tabular}{lrrrrrrr}
\toprule
 & $x_1$ & $x_2$ & $e_1$ & $a_1$ & $a_2$ & RHS & ratio \\
\midrule
row 0 & $0$ & $-10$ & $0$ & $M$ & $4+M$ & $8$ &  \\
row 1 ($x_1$) & $1$ & $-2$ & $0$ & $0$ & $1$ & $2$ & -- \\
row 2 ($e_1$) & $0$ & $-4$ & $1$ & $-1$ & $3$ & $2$ & -- \\
\bottomrule
\end{tabular}
\end{center}

$x_2$ enters (reduced cost $-10$), but its column has no positive entries, so there is no ratio test.

\textbf{Answer:} Unbounded, no optimal solution ($z \to -\infty$). Both artificials have left the basis, so the problem is feasible.

\needspace{12\baselineskip}\section*{Question 2}
\textbf{Standard form}

Minimize: $-z = -2x_1 + 3x_2$

Subject to:
\begin{itemize}
    \item $x_1 + x_2 - e_1 = 3$
    \item $3x_1 + x_2 + s_2 = 6$
    \item $x_1, x_2, e_1, s_2 \ge 0$
\end{itemize}

\textbf{Canonical primal form}

The $\le$ constraint is multiplied by $-1$ so that every constraint is $\ge$.

Minimize: $-z = -2x_1 + 3x_2$

Subject to:
\begin{itemize}
    \item $x_1 + x_2 \ge 3$
    \item $-3x_1 - x_2 \ge -6$
    \item $x_1, x_2 \ge 0$
\end{itemize}

\textbf{Dual form}

The canonical primal is minimize $c^Tx$ subject to $Ax \ge b$, $x \ge 0$, with
\[
c = \begin{bmatrix} -2 \\ 3 \end{bmatrix}, \quad
A = \begin{bmatrix} 1 & 1 \\ -3 & -1 \end{bmatrix}, \quad
b = \begin{bmatrix} 3 \\ -6 \end{bmatrix}.
\]
Its dual is maximize $w = b^Ty$ subject to $A^Ty \le c$, $y \ge 0$, with one $y_i$ for each primal constraint and
\[
A^T = \begin{bmatrix} 1 & -3 \\ 1 & -1 \end{bmatrix}.
\]

Maximize: $w = 3y_1 - 6y_2$

Subject to:
\begin{itemize}
    \item $y_1 - 3y_2 \le -2$ ($x_1$ column)
    \item $y_1 - y_2 \le 3$ ($x_2$ column)
    \item $y_1, y_2 \ge 0$
\end{itemize}

\needspace{12\baselineskip}\section*{Question 3}
The primal is minimize $c^Tx$ subject to $Ax = b$, $x \ge 0$, with
\[
c = \begin{bmatrix} -2 \\ 4 \\ -3 \end{bmatrix}, \quad
A = \begin{bmatrix} 9 & -2 & -8 \\ 3 & 3 & 3 \\ 7 & -5 & 2 \end{bmatrix}, \quad
b = \begin{bmatrix} 5 \\ 7 \\ 9 \end{bmatrix}.
\]
Each equality can be written as a pair of $\ge$ constraints, $Ax \ge b$ and $-Ax \ge -b$, with dual variables $u, v \ge 0$. They only appear as $y = u - v$, so each $y_i$ is free. The dual is maximize $w = b^Ty$ subject to $A^Ty \le c$, $y$ free, with
\[
A^T = \begin{bmatrix} 9 & 3 & 7 \\ -2 & 3 & -5 \\ -8 & 3 & 2 \end{bmatrix}.
\]

Maximize: $w = 5y_1 + 7y_2 + 9y_3$

Subject to:
\begin{itemize}
    \item $9y_1 + 3y_2 + 7y_3 \le -2$ ($x_1$ column)
    \item $-2y_1 + 3y_2 - 5y_3 \le 4$ ($x_2$ column)
    \item $-8y_1 + 3y_2 + 2y_3 \le -3$ ($x_3$ column)
    \item $y_1, y_2, y_3$ free (unrestricted in sign)
\end{itemize}

\needspace{12\baselineskip}\section*{Question 4}
Minimize: $z = -x_1 + 2x_2 - x_3$

Subject to:
\begin{itemize}
    \item $x_1 + 2x_2 + x_3 + s_1 = 12$
    \item $2x_1 + x_2 - x_3 + s_2 = 6$
    \item $-x_1 + 3x_2 + s_3 = 9$
    \item $x_1, x_2, x_3, s_1, s_2, s_3 \ge 0$
\end{itemize}

\needspace{10\baselineskip}\textbf{Iteration 0}

\begin{center}
\begin{tabular}{lrrrrrrrr}
\toprule
 & $x_1$ & $x_2$ & $x_3$ & $s_1$ & $s_2$ & $s_3$ & RHS & ratio \\
\midrule
row 0 & $-1$ & $2$ & $-1$ & $0$ & $0$ & $0$ & $0$ &  \\
row 1 ($s_1$) & $1$ & $2$ & $1$ & $1$ & $0$ & $0$ & $12$ & $12$ \\
row 2 ($s_2$) & $\boxed{2}$ & $1$ & $-1$ & $0$ & $1$ & $0$ & $6$ & $3$ \\
row 3 ($s_3$) & $-1$ & $3$ & $0$ & $0$ & $0$ & $1$ & $9$ & -- \\
\bottomrule
\end{tabular}
\end{center}

$x_1$ and $x_3$ tie at $-1$; $x_1$ enters, $s_2$ leaves.

\needspace{10\baselineskip}\textbf{Iteration 1}

\begin{center}
\begin{tabular}{lrrrrrrrr}
\toprule
 & $x_1$ & $x_2$ & $x_3$ & $s_1$ & $s_2$ & $s_3$ & RHS & ratio \\
\midrule
row 0 & $0$ & $2.5$ & $-1.5$ & $0$ & $0.5$ & $0$ & $3$ &  \\
row 1 ($s_1$) & $0$ & $1.5$ & $\boxed{1.5}$ & $1$ & $-0.5$ & $0$ & $9$ & $6$ \\
row 2 ($x_1$) & $1$ & $0.5$ & $-0.5$ & $0$ & $0.5$ & $0$ & $3$ & -- \\
row 3 ($s_3$) & $0$ & $3.5$ & $-0.5$ & $0$ & $0.5$ & $1$ & $12$ & -- \\
\bottomrule
\end{tabular}
\end{center}

$x_3$ enters, $s_1$ leaves.

\needspace{10\baselineskip}\textbf{Iteration 2}

\begin{center}
\begin{tabular}{lrrrrrrr}
\toprule
 & $x_1$ & $x_2$ & $x_3$ & $s_1$ & $s_2$ & $s_3$ & RHS \\
\midrule
row 0 & $0$ & $4$ & $0$ & $1$ & $0$ & $0$ & $12$ \\
row 1 ($x_3$) & $0$ & $1$ & $1$ & $2/3$ & $-1/3$ & $0$ & $6$ \\
row 2 ($x_1$) & $1$ & $1$ & $0$ & $1/3$ & $1/3$ & $0$ & $6$ \\
row 3 ($s_3$) & $0$ & $4$ & $0$ & $1/3$ & $1/3$ & $1$ & $15$ \\
\bottomrule
\end{tabular}
\end{center}

Row 0 is all $\ge 0$, so this is optimal.

\textbf{Answer:} $x_1 = 6$, $x_2 = 0$, $x_3 = 6$, $z = -12$

\needspace{12\baselineskip}\textbf{Dual form}

Each $\le$ constraint is multiplied by $-1$ to get the canonical form, minimize $c^Tx$ subject to $Ax \ge b$, $x \ge 0$, with
\[
c = \begin{bmatrix} -1 \\ 2 \\ -1 \end{bmatrix}, \quad
A = \begin{bmatrix} -1 & -2 & -1 \\ -2 & -1 & 1 \\ 1 & -3 & 0 \end{bmatrix}, \quad
b = \begin{bmatrix} -12 \\ -6 \\ -9 \end{bmatrix}.
\]
The dual is maximize $w = b^Ty$ subject to $A^Ty \le c$, $y \ge 0$, with
\[
A^T = \begin{bmatrix} -1 & -2 & 1 \\ -2 & -1 & -3 \\ -1 & 1 & 0 \end{bmatrix}.
\]

Maximize: $w = -12y_1 - 6y_2 - 9y_3$

Subject to:
\begin{itemize}
    \item $-y_1 - 2y_2 + y_3 \le -1$ ($x_1$ column)
    \item $-2y_1 - y_2 - 3y_3 \le 2$ ($x_2$ column)
    \item $-y_1 + y_2 \le -1$ ($x_3$ column)
    \item $y_1, y_2, y_3 \ge 0$
\end{itemize}

\needspace{12\baselineskip}\textbf{Complementary slackness}

If a primal variable is positive, its dual constraint is tight. If a primal constraint has slack, its dual variable is 0. The primal optimum has $x_1 = 6$, $x_2 = 0$, $x_3 = 6$, $s_1 = 0$, $s_2 = 0$, $s_3 = 15$.

\begin{center}
\begin{tabular}{ll}
\toprule
Primal & Dual \\
\midrule
$s_3 = 15 > 0$ & $y_3 = 0$ \\
$x_1 = 6 > 0$ & $-y_1 - 2y_2 + y_3 = -1$ \\
$x_3 = 6 > 0$ & $-y_1 + y_2 = -1$ \\
$x_2 = 0$ & no information \\
\bottomrule
\end{tabular}
\end{center}

With $y_3 = 0$, subtracting the third equation from the second gives $-3y_2 = 0$, so $y_2 = 0$ and $y_1 = 1$.

Check: the $x_2$ constraint gives $-2(1) - 0 - 0 = -2 \le 2$, and $y \ge 0$, so $y$ is dual feasible. Also $w = -12(1) = -12 = z^*$. This matches row 0 of the final tableau under $s_1, s_2, s_3$.

\textbf{Answer:} $y_1 = 1$, $y_2 = 0$, $y_3 = 0$, $w = -12$

\section*{Question 5}
Dual problems ask, ``What is each resource worth?'' In HW3 we used optimization to decide which of 3 factories to close. If each factory's capacity is a constraint, its dual value (shadow price) is how much profit the company loses per unit of that capacity. A factory with slack has a dual value of 0. It is not a bottleneck, so it is the best candidate to close. A factory with a high dual value is limiting profit and should stay. Shadow prices only hold for small changes, and closing a factory removes all of its capacity, so we re-solve the model without that factory to confirm. The dual tells us where to look, and re-solving gives the answer.

\end{document}
